% Fichier complété par tools/complete_missing_corrections.py.
% Source : PERF/PERF-06-Precision/73_Bassin/73_Bassin.tex
% Statut : rédigé (6 question(s)).
\begin{corrigebox}[Corrigé rédigé]
\CorrigeQuestion{1}
Pour une boucle ouverte $L_1(p)=C_1(p)G(p)$, le théorème de la valeur finale
        donne $E_S=V_0\lim_{p\to0}1/(1+L_1(p))$ et
        $E_T=\gamma_0\lim_{p\to0}1/[p(1+L_1(p))]$. En remplaçant
        $C_1$, $K_N$, $C$ et $T_m$, on obtient directement les expressions littérales;
        le nombre d'intégrateurs de $L_1$ fixe si ces erreurs sont nulles ou finies.
\CorrigeQuestion{2}
L'exigence impose $|E|\le E_{max}$. On isole donc $C$ dans l'expression
        précédente : $C\ge C_\varepsilon$. Si l'erreur statique est nulle grâce à un
        intégrateur, la condition provient de l'erreur de traînage.
\CorrigeQuestion{3}
En factorisant le numérateur et le dénominateur du correcteur, on identifie
        $C_2(p)=K(1+Tp)^2/p$. L'identification coefficient par coefficient fournit
        $T=T_e$ et $K=C/T_e$ (ou les expressions équivalentes selon la normalisation
        exacte donnée dans l'énoncé).
\CorrigeQuestion{4}
La FTBO corrigée est le produit de tous les blocs de la chaîne directe :
        $L_2(p)=C_2(p)K_NG_m(p)G_{proc}(p)$; on simplifie les facteurs communs et on
        conserve explicitement le nombre d'intégrateurs.
\CorrigeQuestion{5}
Avec $L_2$, $E_S=V_0\lim_{p\to0}1/(1+L_2)$ et
        $E_T=\gamma_0\lim_{p\to0}1/[p(1+L_2)]$. La présence de l'intégrateur annule
        l'erreur d'échelon; le gain de vitesse $K_v=\lim_{p\to0}pL_2(p)$ donne
        $E_T=\gamma_0/K_v$.
\CorrigeQuestion{6}
La condition est donc $K_v\ge\gamma_0/E_{T,max}$. En remplaçant $K_v$ par
        son expression en fonction de $K$, on obtient la borne inférieure recherchée pour
        le gain du correcteur.
\end{corrigebox}
